\chapter{第一族与第二族：容易交出电子的元素}

\begin{kaichang}
去药箱里翻一翻，很可能找到一瓶钙片。看看成分表：碳酸钙，也许还加了点维生素 D。它安安静静地躺在瓶子里，你可以直接嚼碎吞下去，什么都不会发生——这正是它该做的。

现在想象另一幅画面（你只能在网上看视频，或者在课堂上看老师演示，\textbf{绝对不许自己试}）：从煤油里夹出一小块金属钠，用滤纸吸干，丢进一杯水。它立刻熔成一颗银亮的小球，在水面上嘶嘶作响、四处乱窜，有时“砰”的一声燃起黄色火焰，甚至把水杯炸裂。

钠和钙，在周期表上紧挨着，都属于“容易交出电子”的元素。为什么一个暴烈到要泡在油里关押，一个却温顺得能进药瓶？你吞下去的钙片，和金属钙是一回事吗？

先猜一个答案。这一章结束时，我们回来对。
\end{kaichang}

\section{视频里的三兄弟}

请先看一段公开的演示视频（关键词：“碱金属与水反应”，只看正规科普频道或学校课程录像）：三只烧杯，各装半杯水，依次投入米粒大小的锂、钠、钾。你会看到：

\begin{itemize}[itemsep=2pt]
\item \textbf{锂}浮在水面，缓缓移动，嘶嘶冒泡，像一块正在融化的冰；
\item \textbf{钠}反应快得多，放出的热把自己熔成小球（钠的熔点只有 $97.8\,^{\circ}\mathrm{C}$），在水面乱窜，冒泡声尖锐；
\item \textbf{钾}刚一沾水就燃起淡紫色火焰，常常伴随着轻微的爆裂声。
\end{itemize}

再往下看几代：铷和铯一接触水面就剧烈爆炸，所以正规演示里也越来越少见。铷、铯极其昂贵又危险，你在生活中几乎碰不到它们。

这些金属在动手之前是什么样？它们被浸泡在煤油或石蜡油里，像罐头里的水果。取出来的时候，表面灰扑扑的；用小刀切一刀（它们软得像硬奶酪），断面露出银白色的金属光泽——但这光泽只维持几秒钟，很快就暗淡下去，像切开的苹果在空气里变色。

为什么要泡在油里？因为空气对它们来说就是“一锅汤”：氧气会抢它们的电子（断面变暗，就是表面生成了氧化物），水汽更不用提。要关押这些好动分子，必须找一种不含氧、不含水、又不跟它们反应的液体把它们与空气隔开。煤油恰好胜任，还附带一个方便的性质：钠的密度约 \ud{0.97}{g/cm^{3}}，煤油只有 \ud{0.8}{g/cm^{3}} 左右，钠块会沉到瓶底，被油彻底封住——想露头都不行。锂却是个例外：它密度只有 \ud{0.53}{g/cm^{3}}，比煤油还轻，会浮在油面上接触空气，所以锂要封存在固体石蜡里。你看，连“怎么保存”这种问题，答案也写在元素的物理性质里。

同一族的铍、镁、钙、锶、钡（第二族，叫“碱土金属”）脾气缓和一些，但也是“活泼分子”：老师演示的镁条燃烧，用坩埚钳夹住一小段银灰色镁条，点燃后发出刺眼到不能直视的白光，像一小截落进教室的闪电，几十秒后只剩一撮白色粉末，轻得一口气就能吹散。这道白光曾是早期闪光灯和照明弹的光源。金属钙丢进水里，虽不如钠那样暴躁，也会稳定地冒出一串串气泡，水慢慢变浑浊——生成的氢氧化钙溶解度不大，把清水搅成了淡淡的“石灰乳”。

而同一时刻，你家的鸡蛋壳、粉笔、大理石窗台、钙片——这些含第二族元素的东西——安静地待在原地，泡在水里也毫无动静。

这就是本章要拆开的谜：\textbf{同样是这两族的元素，为什么有时候暴烈，有时候温顺？差别到底出在哪里？}

\section{价电子：原子伸在最外面的那只手}

第 9 章讲过电子住“楼层”，第 10 章讲过“最外层电子数决定性格”。这一族的性格，就是从这个规则里长出来的。

锂、钠、钾，最外层都只有 \textbf{1 个}电子；铍、镁、钙，最外层都有 \textbf{2 个}电子。这 1～2 个电子住在最高、最外层的“楼层”上，离原子核最远，还被内层电子挡着（屏蔽效应，第 11 章细讲），是原子全身被“抓得最松”的电子。

现在让它们遇见水。水分子不是省油的灯——它会拉扯金属表面的电子（第 23 章会讲水为什么这么大的本事）。对钠来说，账是这样算的：丢掉那 1 个外层电子要花一点“赎金”（能量），但丢完之后，带正电的钠离子被水分子团团围住，体系放出的能量更多，里外里是赚的；而且反应产物之一氢气还会带着热量跑掉。只要总账是能量下降的，变化就可能发生——原子没有“愿望”，它只是顺着能量的坡往下滚。

丢掉电子之后，原子变成了什么？变成了\textbf{离子}——带电的原子。钠原子丢 1 个电子，变成带 1 个单位正电荷的钠离子；镁原子丢 2 个，变成带 2 个单位正电荷的镁离子。这些离子的外层恰好是住满的“楼层”，能量低，很安稳——这正是“它们为什么偏要丢 1 个或 2 个、而不是丢 3 个”的原因。

\begin{qiaoliang}
凭什么说“只丢 1 个”？证据藏在\textbf{电离能}里——把气态原子的电子一个一个剥下来所需的能量。以钠为例（每摩尔原子计）：
\begin{itemize}[itemsep=1pt]
\item 剥第 1 个电子：约 \ud{496}{kJ/mol}；
\item 剥第 2 个电子：约 \ud{4562}{kJ/mol}——贵了九倍多；
\item 剥第 3 个：约 \ud{6912}{kJ/mol}，依旧昂贵。
\end{itemize}
剥第 1 个电子“便宜”，第 2 个突然涨价近一个数量级——因为第 1 个电子剥掉后，再往下剥就要动住满的第二层了，那层电子离核近、抓得牢。镁的数据同样说明问题：前两个约 \ud{738}{} 和 \ud{1451}{kJ/mol}，第 3 个跳到约 \ud{7733}{kJ/mol}。这个“价格悬崖”的位置，精确地告诉我们每种原子“愿意”丢几个电子：第一族丢 1 个，第二族丢 2 个，各停在自己的悬崖边上。这就是为什么常见离子的电荷可以\textbf{从族号直接读出来}：第一族 $+1$，第二族 $+2$。
\end{qiaoliang}

\section{为什么一个比一个暴躁}

锂只是打转，钠熔成小球，钾直接起火——同一个模子刻出来的脾气，为什么一代比一代烈？

第 11 章讲过这场拉锯战，这里只取结论：同一族从上往下，最外层电子住的“楼层”越来越高，离核越来越远，中间隔的内层电子“棉被”越来越厚，原子核对它的吸引力越来越弱。吸得松，就丢得快：锂的那颗电子要费力一些才被拉走，钾的几乎是一碰就掉。于是和水反应的启动门槛一代比一代低，反应一代比一代猛。

注意一个容易滑过去的区别：\textbf{“容易丢电子”说的是倾向（活泼性），“反应猛不猛”说的是速度（剧烈程度）}，这是两回事，不能画等号。倾向由能量账决定（第 27、28 章细讲），速度由反应条件决定（第 29 章细讲）。三个条件特别要紧：

\begin{itemize}[itemsep=2pt]
\item \textbf{温度}：常温下镁和冷水几乎没动静，放进热水里就稳定冒泡；点燃的镁条才能在空气里烧起来——高温让粒子碰撞更有劲，跨过反应的能量门槛。
\item \textbf{表面积}：整块的钙放进水里只是冒泡，钙粉却可能猛烈反应；打磨成粉的金属，几乎全身都是“前线”。
\item \textbf{表面那层“壳”}：很多活泼金属表面有一薄层氧化物，像穿了雨衣，把水挡在外面。钠切开后迅速变暗，就是在“现场缝制”这件雨衣；铝（第十三族，第 21 章再细讲）比镁还活泼，却能做易拉罐和飞机，全靠它那层致密又结实的氧化膜。
\end{itemize}

所以“活泼”不等于“见什么都炸”。把一颗光滑的铝钉和一杯水放在一起，什么也不会发生——不是它不想丢电子，是雨衣先挡住了。记住这个区别，后面读金属腐蚀、电池时都会用上。

\section{该符号登场了}

现象、粒子、规则都齐了，现在才轮到化学方程式。它们只是把前面讲过的账本压缩成一行字。

钠与水（第一族通用，把 \chem{Na} 换成 \chem{Li}、\chem{K} 即可）：
\[\chem{2Na} + \chem{2H_2O} \rightarrow \chem{2NaOH} + \chem{H_2}\uparrow\]
读法：每 2 个钠原子，各把 1 个电子“转交”给水分子，自己变成钠离子；水分子被拆出氢气，剩下的部分与钠离子组成氢氧化钠（一种强碱，第 32 章细讲），溶解在水里。演示视频里那杯反应后的水，用手摸是烫的，滴上指示剂会显示强碱性——肥皂般的滑腻感（\textbf{不许用手试}，氢氧化钠会腐蚀皮肤）。生成的氢气怎么验明正身？课堂演示有个经典办法：把反应产生的气泡收集起来（或用蘸了肥皂液的导管吹出泡泡），拿燃着的木条一碰，肥皂泡“啪”地一声轻爆——那是氢气与空气混合后燃烧的声音，化学家叫它“爆鸣”。证据链就此闭合：冒出来的确实是氢气，方程式右边的 \chem{H_2} 不是猜的。

镁在空气中燃烧：
\[\chem{2Mg} + \chem{O_2} \rightarrow \chem{2MgO}\]
那道刺眼白光的能量，就来自镁原子和氧原子结成氧化镁时放出的能量（成键放能，第 17、27 章细讲）。烟花和照明弹里掺镁粉，要的就是这道白光；掺锶盐发红光，掺钙盐发橙红光——第二族元素在焰色上也“认亲”。

钙与水（第二族的温和版本）：
\[\chem{Ca} + \chem{2H_2O} \rightarrow \chem{Ca(OH)_2} + \chem{H_2}\uparrow\]

一个数量级估算，感受这些反应的能量。钠与水的反应，每摩尔钠大约放热 \ud{180}{kJ}。演示里常用的一小块钠约 \ud{2.3}{g}，即 0.1 摩尔，放热约 \ud{18}{kJ}。这些热量若全被 \ud{50}{g} 水吸收，按每克水升温 $1\,^{\circ}\mathrm{C}$ 需 \ud{4.2}{J} 计算，足以让水温上升约 $85\,^{\circ}\mathrm{C}$——米大小的金属，几秒钟内把半杯水从室温推到接近沸腾，还顺手点着了氢气。现在你知道为什么钾会直接起火了。

\section{戴维的电铲：从“不可再分”里挖出两种新元素}

\begin{zhengju}
19 世纪初，化学界有一份“不可再分物质”清单——也就是当时认为的元素。清单上，苏打（碳酸钠）和钾碱（碳酸钾）被写得清清楚楚。可拉瓦锡早就嘀咕过：这两种东西的脾气太像了，像是一对兄弟，会不会根本不是“基本物质”，而是某种未知金属的化合物？

怀疑归怀疑，谁也拆不开它们。最猛的火、最烈的酸都试过，苏打和钾碱纹丝不动。直到 1807 年，英国皇家学会的汉弗里·戴维拿到了一件新武器：当时世界上最大的伏打电池组——几百对铜锌片连起来的“电铲”。

他的想法很直接：化学反应分不开的，用电硬拽。电流就是一种“抢电子的力量”，既然这些物质可能是“金属＋其他东西”靠电子转移结合的，那就用更强的电子洪流把它拆开（这就是电解，第 33、36 章细讲）。先电解水溶液，失败了——出来的只是氢气和氧气，水比目标物质更容易被拆。戴维改换思路：把钾碱加热熔化，不让水捣乱，再通电。1807 年 10 月 6 日，阴极上出现了银白色的小液滴，有的立即“啪”地燃烧起来。他记下那种狂喜——据说高兴得在实验室里跳了起来，把记录本都划破了。他给新元素取名 potassium（钾）。几天后，同样的方法拆开了苏打，得到 sodium（钠）。

这个故事里最值得带走的不是“戴维运气好”，而是那条方法论：\textbf{一种物质是不是“元素”，不取决于它多顽固，而取决于你手里有没有更强的拆解工具}。清单可以被新工具改写。反过来，新元素遇空气就着、遇水就炸的暴烈脾气，恰恰解释了为什么它们藏得这么深——在自然界里，它们从来都等不到以金属单质的形式存在，早就把电子交出去、变成安分的离子，锁进矿石和盐湖里了。我们吃下的每一粒盐、每一块石头里都有它们的身影，可“钠”这种金属，地球上没有一块是天然的。
\end{zhengju}

\section{这套模型什么时候会失灵}

“族号读电荷、越下越活泼”是一张很好用的地图，但地图不是领土。它的边界至少有三处。

第一，它描述的是\textbf{单质原子的丢电子倾向}，对已经变成离子的元素完全失效。钠离子 \chem{Na^+} 泡在水里几千年也不闹，因为它已经把电子丢完了，账已经结清——“钠很活泼”这句话的主语永远是金属钠，不是食盐里的钠离子。

第二，实际反应还受那三件“反应条件”的摆布。按活泼性排序，铝在铁前面，可铝锅天天装水炖汤，铁钉却悄悄生锈——氧化膜这层变量，能让纸面排序在现实里翻盘。预测“会不会反应”之前，先问“表面有没有壳、温度多少、接触面积多大”。

第三，钠遇水“爆炸”的细节，比你想象的更微妙。教科书说法是反应放热点燃氢气，氢气爆燃——这对小颗粒基本成立。但高速摄影显示，钠球入水的最初一瞬间，表面会炸出无数尖刺，像蓝色的小刺猬：有研究认为是表面电子瞬间溶进水里，金属表面带上强正电，同号电荷互相排斥把金属“炸开”（称为库仑爆炸），这一步可能比氢气燃烧更早。机制的细节至今仍在讨论——科学在“钠与水”这种百年老反应上都没有写完最后一页。

第四，同一族内部也有“个性”。铍排在第二族最顶端，照理最活泼？恰恰相反——铍表面的氧化膜致密又牢固，把它裹得死死的，金属铍在空气里相当安定；而铍的化合物毒性很强，吸入铍粉尘会严重损伤肺。钡恰好相反，活泼、化合物多可溶、可溶钡盐剧毒。族是“性格倾向”的地图，不是每个成员的简历——具体元素的具体行为，仍要看它所处的形态、环境和剂量。

\begin{wuqu}
“吃钙片就是吃金属钙，补钙补的是钙这种活泼金属。”——错得很有代表性。钙片的成分是碳酸钙（或柠檬酸钙等），里面没有一丝一毫的金属钙，只有\textbf{钙离子} \chem{Ca^{2+}}：钙原子早已把 2 个电子交出去，账结清了，脾气也没了。真把金属钙吞下去，它会在你的唾液和胃液里冒氢气、生成强碱，灼伤消化道——那不是补钙，是服毒。元素、单质、离子是三回事（第 2 章就打过这个预防针）：说“人体需要钙”，指的是钙离子；说“钙遇水冒泡”，指的是金属钙。药瓶里的那位和泡煤油的那位，只是身份证相同，性格早已分道扬镳。
\end{wuqu}

\section{回到那片钙片}

现在可以回答开场的问题了。

钙片为什么温顺？因为它里面的钙早就不是金属钙，而是\textbf{已经丢完电子的钙离子} \chem{Ca^{2+}}，和碳酸根 \chem{CO_3^{2-}} 正负相吸，整整齐齐地锁在碳酸钙的晶体里（离子晶体怎样搭起来，第 18、24 章细讲）。电子的账结清了，能量的坡滚到底了，它没有任何“动力”再跟水或空气折腾。暴烈的是“还没丢电子”的金属单质，温顺的是“已经丢完”的离子——同一元素，两种人生。

你吞进肚里的钙离子，也不会一直游离着：它们中的大多数会被送进骨骼这座“钙矿”砌起来，少数留在血液里执行信号任务，多余的随尿液排出，再由下一顿饭补上——一进一出，维持着精密的平衡。开场的悬念至此解开：药瓶与煤油瓶之间隔着的，不是两种元素，而是同一张身份证的两种电子状态。

这也顺便回答了另一个问题：自然界里为什么找不到金属钠、金属钙？因为它们的单质太“坐不住”，一形成就把电子交给了周围的氧、氯、水，变成离子藏进岩盐、石灰岩、盐湖里。我们今天能见到的所有第一、二族元素，几乎全是“离子形态”；想看到金属单质，必须像戴维那样，用电把电子硬塞回去——然后赶紧泡进煤油，看住它。

\begin{shiyan}
\textbf{鸡蛋壳里的第二族元素}（安全，可在家做）

材料：一个鸡蛋壳（或一小截粉笔，成分都是碳酸钙）、白醋（约 \ud{5}{\%} 的醋酸溶液）、一只透明玻璃杯。

步骤：把鸡蛋壳掰成两三块放进杯子，倒入白醋没过蛋壳，静置观察，过十分钟再看一次，几个小时后再看一次。

观察点：①蛋壳表面立刻出现什么？气泡从哪儿长出来、往哪儿跑？②几小时后蛋壳变软了吗、变薄了吗？③把杯子放一夜，第二天蛋壳还剩什么？

你看到的是：碳酸钙里的钙离子“换了个搭档”——醋酸把碳酸根拆走，放出二氧化碳气泡（就是气泡的真身），钙离子则溶进醋里。啃不动的蛋壳，就这样被弱酸“拆楼”。安全提示：醋酸虽弱也别溅进眼睛；实验后的液体直接倒进下水道冲掉即可。想一想：胃酸比醋酸强，这就是为什么钙片嚼碎吞下去能被溶解吸收——以及为什么贝壳、珊瑚怕海水变酸（第 34 章会回到这件事）。
\end{shiyan}

\section{交了电子之后，它们撑起了你的生活}

第一、二族元素一旦交出电子变成离子，就从“危险分子”变成了生活里随处可见的幕后角色。挑四个看看。

\textbf{骨头是一整座钙矿}。一个成年人体内约有 \ud{1}{kg} 钙，按每摩尔 \ud{40}{g} 估算，约合 25 摩尔、$1.5\times 10^{25}$ 个钙原子，其中约九成九以磷酸钙盐的形式砌在骨骼和牙齿里，既当结构材料又当“钙仓库”。剩下的百分之一溶在血液和细胞液里，却管着肌肉收缩、神经信号这样的大事——血钙浓度的微小波动都要被严格调节（第 62 章讲激素调节时再见）。所以补钙不是“多吃钙就行”：剂量、吸收率、个体差异都有讲究，需不需要补、补多少，该听医生的，而不是广告的。

\textbf{一顿“喝钡”的体检}。做消化道 X 光检查时，医生会让你喝一杯白色的“钡餐”——硫酸钡 \chem{BaSO_4} 悬浊液。等等，钡离子明明有毒，怎么敢喝？关键在溶解度：硫酸钡几乎不溶于水，喝下去穿肠而过，几乎不产生游离的钡离子，却能强烈吸收 X 射线，把肠胃的轮廓涂白给医生看。同一种元素，可溶的钡盐是毒，难溶的钡盐是药——毒性、药效从来都要连着“它在身体里以什么形态、多大剂量出现”一起看。

\textbf{钾是植物的口粮}。庄稼对钾离子的需求量大到“钾肥”与氮肥、磷肥并称肥料三要素：钾离子在植物细胞里调节水分进出和酶的活性。你家香蕉、土豆里的钾，就是这么从土壤里一路走上餐桌的。草木灰（含碳酸钾）是历史最悠久的钾肥——名字里的“钾”（kalium）正是从“草木灰”来的。

\textbf{最轻的金属驮着最重的电量}。锂是最轻的金属，也是最“舍得”丢电子的元素之一：同样一克材料，锂能交出的电子多、交得干脆，于是成了充电电池的明星（第 36 章讲电池时细拆）。你的手机、电动车里那块锂电池，本质上是把“锂丢电子”这个本章主角脾气，装进可控的笼子里慢慢用。锂盐还是治疗双相情感障碍的老牌药物——剂量窗口很窄，必须在医生监测下使用，这又是一个“形态与剂量决定一切”的例子。

还有一条暗线别忘了：钠和钾丢出去的电子，总得有个“下家”接手。在水里，下家是水分子，产物是氢气。但如果让钠遇见一个\textbf{特别爱抢电子}的对手呢？那将是另一种更彻底、更常见的结合方式——食盐就是这样来的。周期表最右边那一族，住着一群“见电子就抢”的元素，和本章这群“见电子就让”的元素恰好相反。下一章，我们去拜访它们。

\begin{tiaozhan}
“活泼性”其实是\textbf{可以量出来}的，不用只靠看视频比谁的动静大。方法是把金属放进它自己离子的溶液里，接上电路，测它“推电子”的劲头有多足——这就是电极电势（第 36 章细讲）。测出来的一串数字（单位伏特，数值越负代表越舍得丢电子）：钾约 $-2.93$、钙约 $-2.87$、钠约 $-2.71$、镁约 $-2.37$、锂约 $-3.04$。注意两个“反直觉”：第一，按这串数字，锂才是全表最舍得丢电子的，可是演示视频里锂遇水最温和——倾向最大不代表速度最快，锂熔点高（$180.5\,^{\circ}\mathrm{C}$），放出的热不足以把它熔成乱跑的小球，接触面积小，速度就上不去；倾向与速度再次分家。第二，钙比钠数字更负，遇水却比钠温和，原因之一是生成的氢氧化钙溶解度小，会糊在钙块表面当“锅盖”。用仪器代替“看热闹”，用机制解释“意外”——这是化学家处理排序问题的标准姿势。跳过本节不影响主线。
\end{tiaozhan}

\section*{练一练}

\begin{sikaoti}
\item 参照本章“电离能价格悬崖”，不看资料预测：铝（第十三族，最外层 3 个电子）的第几电离能会出现大跳跃？它的常见离子带什么电荷？画出你的推理。
\item 把金属钙分别放进冷水和稀盐酸，预测哪个反应更剧烈，并说明理由（提示：考虑“抢电子的对手”强弱）。再预测：把同样大小的钙块换成钙粉，会发生什么变化？
\item 一位同学说：“钠既然这么活泼，那把食盐撒进水里一定会冒气泡。”找出这句话里至少两处概念混淆，并写一段纠正它的话。
\item 画一幅能量示意图：横轴是“反应进程”，纵轴是“能量”，分别画出“金属钠＋水”与“钠离子＋氯离子（食盐溶于水）”两条线的起点和终点高低，并用一句话说明哪边“坐得住”、为什么。
\item 钡餐用的是几乎不溶的硫酸钡，而不是可溶的氯化钡。用“离子形态与剂量”的逻辑解释：为什么前者安全、后者危险？如果某次检查误用了可溶钡盐，从化学上你能想到什么急救方向（只谈原理）？
\item 锂电池广告说“锂最轻，所以电池容量大”。这句话对了一半。用本章的“丢电子”语言把另一半补上：除了轻，锂还有什么关键脾气让它适合做电池？
\end{sikaoti}
